% Quarantined proof corresponding exactly to statement.tex.
% It is not a live-book input and has not received production verification.

\begin{proof}
We first isolate three elementary points which will be used twice.

\medskip\noindent\emph{Local extended-real convexity on a convex set.}
Let \(H\) be a convex subset of a finite-dimensional affine space and let
\(v:H\to\mathbf R\cup\{-\infty\}\) be locally extended-real convex.  Then
\(v\) is extended-real convex on \(H\).  It is enough to restrict to a
segment and prove the assertion on an interval \(I\).  On
\(I^\circ\), both
\[
 \{t:v(t)\in\mathbf R\}
 \quad\text{and}\quad
 \{t:v(t)=-\infty\}
\]
are open.  For the first set, if a finite point had arbitrarily close
\(-\infty\) points, write the finite point as a strict convex combination
of one of them and a point on its other side, inside one interval of local
convexity.  For the second set, if \(v(t)=-\infty\), then every sufficiently
nearby \(s\) lies strictly between \(t\) and another point of the same
local convexity interval; the convexity inequality forces
\(v(s)=-\infty\).  Connectedness of \(I^\circ\) shows that \(v\) is either
finite throughout \(I^\circ\) or identically \(-\infty\) there.

In the first case an included endpoint cannot have value \(-\infty\):
local convexity there would propagate \(-\infty\) to a one-sided interior
collar.  Thus \(v\) is finite on \(I\), and the usual subdivision argument
applies: on a partition fine enough that each pair of consecutive
subintervals lies in a local convexity interval, the successive secant
slopes are nondecreasing.  Summing the slope inequalities proves global
convexity.  In the second case every convexity inequality at a point of
\(I^\circ\) has left side \(-\infty\), while the endpoint cases are
tautological.  This proves the assertion, including possible finite values
at the two endpoints of an otherwise polar segment.

\medskip\noindent\emph{Continuous convexification on a polytope.}
Let \(K\) be a compact convex polytope in a \(d\)-dimensional affine space
and let \(\phi:K\to\mathbf R\) be continuous.  Put
\begin{equation}
 (\operatorname{co}_K\phi)(x)
 =
 \inf\left\{
       \sum_i\alpha_i\phi(x_i):
       \alpha_i\geq0,\ \sum_i\alpha_i=1,\
       \sum_i\alpha_i x_i=x
      \right\}.
 \label{eq:p071-polyhedral-convex-envelope}
\end{equation}
This is the largest convex function on \(K\) bounded above by \(\phi\).
It is finite and continuous.

For completeness, every combination in
\eqref{eq:p071-polyhedral-convex-envelope} can be replaced, without
increasing its value, by one with at most \(d+1\) nonzero coefficients.
Indeed, if there are more, affine dependence gives numbers \(\beta_i\),
not all zero, with
\[
             \sum_i\beta_i=0,\qquad \sum_i\beta_i x_i=0.
\]
Move the weights along one of the two directions
\(\alpha_i+t\beta_i\) for which the objective does not increase, until a
coefficient becomes zero, and iterate.  It follows that the infimum is
attained on a compact parameter space.  If \(x_j\to x\), choose minimizing
representations for \(x_j\) and pass to a convergent subsequence.  The
limiting representation is admissible for \(x\), which proves
\[
 \operatorname{co}_K\phi(x)
 \leq\liminf_j\operatorname{co}_K\phi(x_j).
\]
Thus the envelope is lower semicontinuous.  Convexity follows by combining
two admissible representations.

It remains only to check upper semicontinuity at the boundary of \(K\).
Fix \(x\in K\).  A finite triangulation of \(K\) may be refined so that
\(x\) is a vertex and the closed star of \(x\) is a neighbourhood of
\(x\) in \(K\).  If \(x_j\to x\), pass through the finitely many simplices
of that star.  In a simplex containing \(x\), the barycentric coefficient
of \(x\) tends to one.  Convexity bounds the envelope at \(x_j\) by the
same barycentric combination of its values at \(x\) and at the other
vertices.  The latter values are finite, hence
\[
 \limsup_j\operatorname{co}_K\phi(x_j)
 \leq\operatorname{co}_K\phi(x).
\]
This proves continuity.  Notice that polyhedrality is used exactly in
this boundary argument.

The convexification operation is nonexpanding in the uniform norm:
\begin{equation}
 \left\|\operatorname{co}_K\phi-\operatorname{co}_K\psi\right\|_K
 \leq \|\phi-\psi\|_K .
 \label{eq:p071-polyhedral-envelope-nonexpanding}
\end{equation}
This follows immediately by applying
\eqref{eq:p071-polyhedral-convex-envelope} to the same barycentric
representations in both directions.

\medskip\noindent\emph{Smooth convex majorants on one fixed polytope.}
If \(a:K\to\mathbf R\) is continuous and convex, then for every
\(\varepsilon>0\) there is a smooth convex function \(g\) on the whole
affine span of \(K\) such that
\begin{equation}
                       a\leq g\leq a+\varepsilon
                       \qquad\text{on }K.
 \label{eq:p071-polyhedral-smooth-majorant}
\end{equation}
To see this, extend \(a\) by \(+\infty\) off \(K\).  The resulting function
is proper, convex, and lower semicontinuous.  Strict separation of its
closed epigraph from the point \((x,a(x)-\delta)\) gives, for every
\(x\in K\) and \(\delta>0\), an affine function \(\ell_x\) such that
\[
                         \ell_x\leq a\text{ on }K,
              \qquad     \ell_x(x)>a(x)-\delta .
\]
The sets on which \(\ell_x>a-\delta\) form an open cover of \(K\).
After taking a finite subcover and putting
\[
                              p=\max_j\ell_{x_j},
\]
we obtain a globally defined convex piecewise-affine function satisfying
\[
                              a-\delta<p\leq a
                              \qquad\text{on }K.
\]
Choose a nonnegative smooth compactly supported probability kernel
\(\rho\) of barycentre zero.  The function \(p\) is globally Lipschitz;
if \(L\) is a Lipschitz constant, then
\[
                  p\leq p*\rho_r\leq p+Lr .
\]
The first inequality is Jensen's inequality.  Taking
\(\delta=\varepsilon/3\), choosing \(r\) with
\(Lr<\varepsilon/3\), and setting
\[
                              g=p*\rho_r+\delta
\]
proves \eqref{eq:p071-polyhedral-smooth-majorant}.  If \(L=0\), any
sufficiently small \(r\) works.  This argument is uniform for any finite
family of continuous convex functions by choosing the smoothing radii
separately.

\medskip\noindent
We now turn to the comparison sheaf.  Its defining condition is local, so
it is a sheaf.  Every smooth psh generator belongs to
\(\mathscr C_X\): after a clopen shrinking it is independent of \(z\);
take \(Y\) to be a point.  Positive semidefiniteness of the Hessian implies
convexity on each convex relative neighbourhood in \(P\).

We prove that \(\mathscr C_X\) is stable under decreasing sequential
limits.  Let \(u_n\downarrow u\) on an open subset \(W\subset X\), with
\(u_n\in\mathscr C_X(W)\), and fix \((z_0,x_0)\in W\).  Local compactness
and the clopen basis give a compact clopen neighbourhood \(C\) of \(z_0\).
Choose bounded open polyhedra \(O_1,O_0\subset A\), both containing
\(x_0\), such that
\[
 \overline{O_1}\subset O_0,\qquad
 C\times\bigl(P\cap\overline{O_0}\bigr)\subset W.
\]
After making the polyhedra smaller if necessary, we may and do suppose that
\[
       V=P\cap O_1\quad\text{is convex},\qquad
       K=P\cap\overline{O_0}\quad\text{is a compact convex polytope}.
 \label{eq:p071-polyhedral-compact-core}
\]

For each fixed \(n\), compactness gives a finite product cover of
\(C\times K\) by sets
\begin{equation}
                           C_{n,a}\times V_{n,a},
 \label{eq:p071-polyhedral-local-factor-cover}
\end{equation}
where \(C_{n,a}\) is compact clopen in \(C\), \(V_{n,a}\) is relatively
open in \(P\), and on this product \(u_n\) factors through a continuous
surjection
\[
                  q_{n,a}:C_{n,a}\twoheadrightarrow Y_{n,a}
\]
onto a compact metrizable zero-dimensional space.  Choose a countable
clopen base \(\mathcal D_{n,a}\) of each \(Y_{n,a}\).  Let
\(\mathcal B\) be the Boolean subalgebra of the clopen subsets of \(C\)
generated by all \(C_{n,a}\) and all
\[
                     q_{n,a}^{-1}(D),
                     \qquad D\in\mathcal D_{n,a}.
\]
There are only countably many generators.  The evaluation map
\begin{equation}
 q:C\longrightarrow
 Y:=q(C)\subseteq\{0,1\}^{\mathcal B}
 \label{eq:p071-polyhedral-common-quotient}
\end{equation}
is a continuous surjection onto a compact metrizable zero-dimensional
space.

Every \(u_n|_{C\times V}\) is constant on the fibres of
\(q\times\operatorname{id}_V\).  Indeed, suppose \(q(z)=q(z')\) and
\(x\in V\).  Choose a member of
\eqref{eq:p071-polyhedral-local-factor-cover} containing \((z,x)\).
The Boolean coordinates first show that \(z'\in C_{n,a}\), and then that
\(q_{n,a}(z)=q_{n,a}(z')\), since the chosen clopen base separates points.
The local factorization gives \(u_n(z,x)=u_n(z',x)\).  The limit therefore
factors as
\begin{equation}
                           u(z,x)=\bar u(q(z),x)
                           \qquad(C\times V).
 \label{eq:p071-polyhedral-limit-factor}
\end{equation}

The function \(u\) is upper semicontinuous: each \(u_n\) is locally upper
semicontinuous and
\[
                         \{u<t\}=\bigcup_n\{u_n<t\}.
\]
Upper semicontinuity descends to \(\bar u\) without any unsupported
product-quotient assertion.  If \(\bar u(y,x)<t\), cover the compact fibre
\(q^{-1}(y)\times\{x\}\) by finitely many product neighbourhoods
\(N_j\times D_j\) contained in \(\{u<t\}\).  Set
\(N=\bigcup_jN_j\) and \(D=\bigcap_jD_j\).  Since \(q\) is closed,
\[
                         Y\setminus q(C\setminus N)
\]
is an open neighbourhood of \(y\), and its product with \(D\) lies in
\(\{\bar u<t\}\).

For fixed \(z\), the restrictions \(x\mapsto u_n(z,x)\) are locally
extended-real convex.  The first elementary point makes them convex on
\(V\).  Passing to the limit in every finite convexity inequality shows
that \(x\mapsto u(z,x)\), and hence
\(x\mapsto\bar u(y,x)\), is extended-real convex.  Thus
\(\mathscr C_X\) is decreasing-sequentially closed.  Since it contains
the generators, minimality gives
\begin{equation}
                         \operatorname{Psh}_X\subseteq\mathscr C_X.
 \label{eq:p071-polyhedral-minimality}
\end{equation}

It remains to prove that every comparison germ is already in the first
stage.  Let \(u\in\mathscr C_X\) have a factorization
\eqref{eq:p071-polyhedral-factorization} on \(C\times V_0\).  Choose a
compact convex polytope \(K\subset V_0\) and a relatively open polyhedral
neighbourhood \(V\) of \(x_0\) such that
\begin{equation}
                         \overline V^{\,P}\subset K\subset V_0.
 \label{eq:p071-polyhedral-nested-core}
\end{equation}
Here the closure is taken in \(P\).  Such a choice is obtained by
intersecting \(P\) with two nested bounded open polyhedral boxes in
\(A\).  We construct every approximant on the one fixed set \(C\times V\).

Give \(Y\times K\) a compatible product metric \(d\).  The upper
semicontinuous function \(\bar u\) is bounded above there; fix an upper
bound \(M\).  For \(k\geq1\), put
\begin{align}
 h_k(s)&=\max\{\bar u(s),-k\},\\
 \Phi_k(s)&=\sup_{t\in Y\times K}
                 \bigl(h_k(t)-k\,d(s,t)\bigr).
 \label{eq:p071-polyhedral-usc-regularization}
\end{align}
Each \(\Phi_k\) is finite, continuous, and \(k\)-Lipschitz.  Moreover,
\begin{equation}
                \Phi_{k+1}\leq\Phi_k,\qquad
                \Phi_k\geq\bar u,\qquad
                \Phi_k\downarrow\bar u .
 \label{eq:p071-polyhedral-usc-convergence}
\end{equation}
Only the last assertion requires a check.  Fix \(s\) and a real number
\(a>\bar u(s)\), with \(a\) arbitrary when \(\bar u(s)=-\infty\).
Choose \(b\) with \(\bar u(s)<b<a\).  Upper semicontinuity gives a metric
ball about \(s\) on which \(\bar u<b\).  For large \(k\), also \(-k<b\).
On this ball the expression under the supremum is at most \(b\); outside a
smaller fixed ball it is at most
\(\max\{M,-1\}\) minus \(k\) times a fixed positive number.  It is therefore
less than \(a\) for large \(k\).  This proves the last assertion.

For \(y\in Y\), convexify only in the polyhedral variable:
\begin{equation}
 c_k(y,x)
 =\operatorname{co}_K\bigl(\Phi_k(y,\cdot)\bigr)(x).
 \label{eq:p071-polyhedral-fibre-convexification}
\end{equation}
Because \(x\mapsto\bar u(y,x)\) is extended-real convex and
\(\bar u\leq\Phi_k\), every barycentric representation of \(x\) gives
\[
       \bar u(y,x)
       \leq\sum_i\alpha_i\bar u(y,x_i)
       \leq\sum_i\alpha_i\Phi_k(y,x_i).
\]
Taking the infimum yields
\begin{equation}
                     \bar u\leq c_k\leq\Phi_k,
        \qquad       c_{k+1}\leq c_k,
        \qquad       c_k\downarrow\bar u .
 \label{eq:p071-polyhedral-convexified-convergence}
\end{equation}
The function \(c_k\) is jointly continuous on \(Y\times K\).  Continuity
in \(x\) is the second elementary point.  For \(y,y'\in Y\),
\eqref{eq:p071-polyhedral-envelope-nonexpanding} and the Lipschitz bound
for \(\Phi_k\) give
\[
 \sup_{x\in K}|c_k(y,x)-c_k(y',x)|
 \leq
 \sup_{x\in K}|\Phi_k(y,x)-\Phi_k(y',x)|
 \leq k\,d_Y(y,y'),
\]
after choosing the product metric in sum form.  Separate continuity with
this uniform estimate gives joint continuity.

Choose nested finite clopen partitions \(\mathcal P_k\) of \(Y\) whose
meshes tend to zero.  Such partitions are obtained by taking finite clopen
covers of successively smaller mesh and refining all preceding covers.
For \(Q\in\mathcal P_k\), set
\begin{equation}
             a_{k,Q}(x)=\max_{y\in Q}c_k(y,x),
             \qquad
             a_k(y,x)=a_{k,Q_k(y)}(x),
 \label{eq:p071-polyhedral-clopen-envelope}
\end{equation}
where \(Q_k(y)\) is the member containing \(y\).  Each \(a_{k,Q}\) is a
finite continuous convex function on \(K\).  Refinement and
\(c_{k+1}\leq c_k\) give
\begin{equation}
                              a_{k+1}\leq a_k.
 \label{eq:p071-polyhedral-clopen-monotonicity}
\end{equation}
Furthermore,
\begin{equation}
                              a_k\downarrow\bar u
                              \quad\text{on }Y\times K.
 \label{eq:p071-polyhedral-clopen-convergence}
\end{equation}
The lower bound follows from \(a_k\geq c_k\geq\bar u\).  For the upper
bound, fix \((y,x)\), choose \(a>\bar u(y,x)\), and then choose
\(b\) strictly between these two numbers.  Choose a metric ball on which
\(\bar u<b\).  For large \(k\), the set \(Q_k(y)\times\{x\}\) lies close
enough to \((y,x)\), and \(-k<b\).  In
\eqref{eq:p071-polyhedral-usc-regularization}, points in the chosen ball
contribute at most \(b\), whereas points outside it pay a fixed positive
distance times \(k\).  Consequently
\[
       \sup_{y'\in Q_k(y)}\Phi_k(y',x)<a
\]
for large \(k\).  Since \(a_{k,Q_k(y)}\leq
\sup_{y'\in Q_k(y)}\Phi_k(y',x)\), this proves
\eqref{eq:p071-polyhedral-clopen-convergence}, also at a
\(-\infty\)-valued point.

Put \(\varepsilon_k=4^{-k}\).  The smooth-majorant construction, applied
to the finite family indexed by \(\mathcal P_k\), gives smooth convex
functions \(g_{k,Q}\) on \(A\) such that
\begin{equation}
                        a_{k,Q}\leq g_{k,Q}
                        \leq a_{k,Q}+\varepsilon_k
                        \qquad\text{on }K.
 \label{eq:p071-polyhedral-family-smoothing}
\end{equation}
Define \(w_k\) on \(C\times V\) by
\begin{equation}
                 w_k(z,x)=g_{k,Q}(x)
                 \quad\text{if }q(z)\in Q\in\mathcal P_k.
 \label{eq:p071-polyhedral-smooth-piecewise}
\end{equation}
The inverse images \(q^{-1}(Q)\) form a finite clopen partition of \(C\)
by subsets which are clopen in \(Z\).  Hence \(w_k\) is smooth for the
stated product structure.  Its Hessian in the \(A\)-directions is positive
semidefinite, so \(w_k\in\operatorname{Psh}_{X,0}(C\times V)\).
Equations \eqref{eq:p071-polyhedral-clopen-monotonicity} and
\eqref{eq:p071-polyhedral-family-smoothing} imply
\begin{equation}
                             w_{k+1}\leq w_k+\varepsilon_{k+1}.
 \label{eq:p071-polyhedral-almost-monotone}
\end{equation}

Set
\[
                 \tau_k=\sum_{j\geq k+1}\varepsilon_j,
                 \qquad G_k=w_k+\tau_k .
\]
Since \(\tau_k=\varepsilon_{k+1}+\tau_{k+1}\),
\eqref{eq:p071-polyhedral-almost-monotone} gives
\(G_{k+1}\leq G_k\).  Finally,
\eqref{eq:p071-polyhedral-clopen-convergence},
\eqref{eq:p071-polyhedral-family-smoothing}, and
\(\tau_k\to0\) show that
\[
                         G_k\downarrow u
                         \qquad\text{on }C\times V.
\]
This is \eqref{eq:p071-polyhedral-fixed-domain}.  We have proved
\[
                         \mathscr C_X
                         \subseteq\operatorname{Psh}_{X,1}.
\]
Together with \eqref{eq:p071-polyhedral-minimality} and the tautological
inclusion
\(\operatorname{Psh}_{X,1}\subseteq\operatorname{Psh}_X\), this proves
\eqref{eq:p071-polyhedral-collapse}.

The proof is local and uses only the generator sheaf.  It therefore
transports through any generator-preserving chart and proves the asserted
chartwise extension.  For a rational convex polyhedron with its induced
Chambert--Loir--Ducros affine tropical structure, smooth functions are
locally restrictions of ambient smooth functions; positivity of
\(\mathrm d'\mathrm d''f\) is positive semidefiniteness of the Hessian on
the affine span.  A locally constant affine structure on \(Z\) becomes
constant after a clopen shrinking, and the product construction gives
exactly the generator sheaf used above.  This verifies the claimed
compatibility while leaving untouched the excluded nonconvex and branched
settings.
\end{proof}

% Proof corresponding exactly to statement.tex in this isolated dossier.
% No production call has been made for this follow-up candidate.

\begin{proof}
We use the following already frozen and production-verified local result,
denoted by \((\mathrm M)\).  If \(Z\) is locally compact Hausdorff with a
clopen basis and \(P\) is any nonempty convex polyhedron, including an
unbounded, lower-dimensional, or zero-dimensional one, then on
\(Z\times P\) the first local decreasing-limit stage of the
single-chamber smooth psh generators, the full decreasing-sequentially
generated psh sheaf, and the Stone-factorized fibrewise-convex comparison
sheaf coincide.  Moreover, every germ has a decreasing approximation by
smooth psh generators on one fixed product
\(C\times W\), where \(C\) is compact clopen and \(W\) is a relatively
open convex polyhedral neighbourhood.  The exact statement, proof, source
chain, hostile audit, and production receipt for \((\mathrm M)\) are pinned
in \texttt{SOURCE\_AND\_OBLIGATION\_MATRIX.md}.

We first record the elementary transport argument needed to apply
\((\mathrm M)\).  Let \(h:T\to T'\) be a homeomorphism, and let
\(\mathcal G_0\) and \(\mathcal G'_0\) be generator sheaves such that
\begin{equation}
                   f\in\mathcal G_0
             \quad\Longleftrightarrow\quad
                   f\circ h^{-1}\in\mathcal G'_0
 \label{eq:p071-complex-stage-zero-transport}
\end{equation}
locally, in both directions.  A local sequence
\(f_j\in\mathcal G_0\) with \(f_j\downarrow f\) transports term by term to
\(f_j\circ h^{-1}\downarrow f\circ h^{-1}\).  Applying the same argument
to \(h^{-1}\) proves that \(h\) identifies the two first local
decreasing-limit stages.

It also identifies the full generated sheaves.  Indeed, the inverse image
under \(h\) of the full generated sheaf on \(T'\) is a sheaf containing
\(\mathcal G_0\) and stable under local decreasing sequential limits.
Minimality gives one inclusion; applying this to \(h^{-1}\) gives the
reverse inclusion.  Thus
\begin{equation}
 h_*(\mathcal G_1)=\mathcal G'_1,
 \qquad
 h_*(\mathcal G)=\mathcal G'.
 \label{eq:p071-complex-all-stage-transport}
\end{equation}
This uses neither a partition of unity nor a countability property of the
underlying spaces.

We will also use that this construction commutes with restriction to an
open subset.  The assertion is immediate for the first local stage.  For
the full sheaf, let \(S\subset T\) be open and write \(\mathcal G_S\) for
the sheaf generated on \(S\) by \(\mathcal G_0|_S\).  One inclusion follows
because \(\mathcal G|_S\) is decreasing-sequentially closed and contains
the restricted generators.  For the reverse inclusion, consider the
subsheaf of \(\mathcal G\) consisting of sections whose germs at points of
\(S\) lie in \(\mathcal G_S\), with no extra condition away from \(S\).
This subsheaf contains \(\mathcal G_0\) and is stable under local decreasing
sequential limits.  The minimality of \(\mathcal G\) therefore forces every
germ of \(\mathcal G|_S\) to lie in \(\mathcal G_S\).  Hence
\begin{equation}
                    \mathcal G|_S=\mathcal G_S.
 \label{eq:p071-complex-restriction}
\end{equation}

Fix \((z_0,p_0)\in X\), and choose a chart
\(\varphi:U\to V\subset P\) from
\eqref{eq:p071-complex-chart}.  On a clopen neighbourhood in \(Z\), the map
\begin{equation}
               \operatorname{id}\times\varphi:
               Z\times U\longrightarrow Z\times V
 \label{eq:p071-complex-product-chart}
\end{equation}
identifies the stage-zero sheaves by hypothesis.  Consequently
\eqref{eq:p071-complex-all-stage-transport} identifies both the first and
the full psh sheaves with their restrictions from the product model
\(Z\times P\), using \eqref{eq:p071-complex-restriction} for the open chart
target \(Z\times V\).

We next check that the comparison condition does not depend on the chart.
Suppose first that a germ has the factorization
\eqref{eq:p071-complex-comparison} in a chart \(\varphi_\alpha\).  By
\((\mathrm M)\), its coordinate representative belongs to the first-stage
model sheaf.  On an overlap, the transition map is generator-preserving;
the first equality in \eqref{eq:p071-complex-all-stage-transport} therefore
places the representative in the first-stage model sheaf of any other
chart \(\varphi_\beta\).  A second application of \((\mathrm M)\), after a
convex shrinking around the image of the point, produces exactly a compact
metrizable zero-dimensional quotient and an upper-semicontinuous
fibrewise-convex factorization in the \(\varphi_\beta\)-coordinates.  The
same argument with \(\alpha\) and \(\beta\) interchanged proves equivalence.
The comparison property is therefore intrinsic.  Since its condition is
local in any chart, it defines a sheaf \(\mathscr C_X\).

Now let \(u\) be a germ of \(\operatorname{Psh}_X\) at
\((z_0,p_0)\).  Transport through
\eqref{eq:p071-complex-product-chart} and
\eqref{eq:p071-complex-all-stage-transport}.  The model theorem
\((\mathrm M)\) places the transported germ in the model comparison sheaf,
and hence in the model first-stage sheaf.  Transporting back proves
\begin{equation}
                  \operatorname{Psh}_X
                  \subseteq\mathscr C_X
                  \subseteq\operatorname{Psh}_{X,1}.
 \label{eq:p071-complex-two-inclusions}
\end{equation}
The tautological inclusion
\(\operatorname{Psh}_{X,1}\subseteq\operatorname{Psh}_X\) completes the
proof of \eqref{eq:p071-complex-collapse}.

The fixed-neighbourhood assertion is equally local but is worth spelling
out.  Apply the fixed-domain clause of \((\mathrm M)\) to the coordinate
representative of \(u\).  After shrinking inside \(V\), it gives a compact
clopen \(C\), a relatively open convex polyhedral neighbourhood \(W\), and
model generators \(g_k\) on the one set \(C\times W\), with
\(g_{k+1}\leq g_k\) and \(g_k\downarrow u\) there.  Set
\begin{equation}
        G_k(z,p)=g_k\bigl(z,\varphi(p)\bigr).
 \label{eq:p071-complex-pulled-approximants}
\end{equation}
The generator-preserving condition makes every \(G_k\) a stage-zero
generator on the same set
\(C\times\varphi^{-1}(W)\), and order and pointwise convergence are
preserved.  This proves \eqref{eq:p071-complex-fixed-star}.  Since
\(\varphi^{-1}(W)\) is one ordinary neighbourhood of \(p_0\), it contains
the portions of every cell incident to the germ.  Thus the model
construction smooths across all chart faces simultaneously.  Separate
chamberwise sequences are never chosen and hence do not have to be glued.

This also explains the loop case.  The proof chooses only one chart at the
given germ.  Continuing that chart around a loop is unnecessary, so the
composition of affine transition maps around the loop does not enter any
step.  Conversely, nothing in the argument makes the local sequences agree
term by term on distinct chart overlaps.  It therefore proves the sheaf
identity, whose first stage is local by definition, but it proves no
single-global-sequence assertion.

We verify the asserted obstruction to weakening the hypotheses.  The
two-ray complex
\[
        L=\mathbf R_{\geq0}e_1\cup\mathbf R_{\geq0}e_2
\]
is face-to-face, and its vertex is incident to exactly two one-dimensional
cells, so its underlying space is a topological one-manifold.  The Hessian
of the function \(F\) in \eqref{eq:p071-complex-v-function} is zero.
Hence its restriction is an ambient-induced cellwise smooth psh generator.
It is continuous and the two restrictions agree at the common face.

Any cellwise affine homeomorphism from a neighbourhood of the vertex to an
interval sends the two rays to opposite half-intervals.  After possibly
interchanging them, it has the form
\[
                 te_1\longmapsto at,
        \qquad  te_2\longmapsto-bt
        \qquad(a,b>0).
\]
In the interval coordinate \(s\), the transported restriction of \(F\)
has slope \(b^{-1}>0\) to the left of zero and slope \(-a^{-1}<0\) to the
right.  A convex function must have nondecreasing one-sided slopes, so this
function is not convex; for \(a=b=1\) it is \(-|s|\).  Therefore no such
unfolding identifies this complete cellwise generator germ with the
single-convex-interval generator germ.  Agreement only on the common
zero-dimensional face does not impose the missing slope inequality.

Finally, all dimensional edge cases are inherited without change from
\((\mathrm M)\).  When \(P\) has dimension zero, its affine span is a
point, fibrewise convexity is vacuous, and the model proof uses no maximum
over an empty set.  For a lower-dimensional \(P\), every Hessian and every
relative interior is taken in \(\operatorname{Aff}(P)\).  At a boundary or
corner, \(W\) is relatively open in \(P\), so the fixed domain remains an
ordinary neighbourhood in \(|\Sigma|\).  Unboundedness causes no problem
because \((\mathrm M)\) constructs the approximation on a bounded compact
core after a local shrinking.  None of these observations supplies a chart
at a branched or singular-link germ, which is why those cases remain outside
the assertion.
\end{proof}
